\documentclass[10pt,a4paper]{amsart}
\usepackage{amsmath,amssymb,amsthm,graphicx}
\usepackage{verbatim}
\newenvironment{reference}{\comment}{\endcomment}
%\newenvironment{reference}{}{}
\newenvironment{slogan}{\comment}{\endcomment}
\newenvironment{history}{\comment}{\endcomment}

% For commutative diagrams we use Xy-pic
\usepackage[all]{xy}

% We use 2cell for 2-commutative diagrams.
\xyoption{2cell}
% Theorem environments.
%
\theoremstyle{plain}
\newtheorem{theorem}[subsection]{Theorem}
\newtheorem{proposition}[subsection]{Proposition}
\newtheorem{lemma}[subsection]{Lemma}

\theoremstyle{definition}
\newtheorem{definition}[subsection]{Definition}
\newtheorem{example}[subsection]{Example}
\newtheorem{exercise}[subsection]{Exercise}
\newtheorem{situation}[subsection]{Situation}

\theoremstyle{remark}
\newtheorem{remark}[subsection]{Remark}
\newtheorem{remarks}[subsection]{Remarks}

\numberwithin{equation}{subsection}

% Macros
%
\def\lim{\mathop{\mathrm{lim}}\nolimits}
\def\colim{\mathop{\mathrm{colim}}\nolimits}
\def\Spec{\mathop{\mathrm{Spec}}}
\def\Hom{\mathop{\mathrm{Hom}}\nolimits}
\def\Ext{\mathop{\mathrm{Ext}}\nolimits}
\def\SheafHom{\mathop{\mathcal{H}\!\mathit{om}}\nolimits}
\def\SheafExt{\mathop{\mathcal{E}\!\mathit{xt}}\nolimits}
\def\Sch{\mathit{Sch}}
\def\Mor{\mathop{\mathrm{Mor}}\nolimits}
\def\Ob{\mathop{\mathrm{Ob}}\nolimits}
\def\Sh{\mathop{\mathit{Sh}}\nolimits}
\def\NL{\mathop{N\!L}\nolimits}
\def\CH{\mathop{\mathrm{CH}}\nolimits}
\def\proetale{{pro\text{-}\acute{e}tale}}
\def\etale{{\acute{e}tale}}
\def\QCoh{\mathit{QCoh}}
\def\Ker{\mathop{\mathrm{Ker}}}
\def\Im{\mathop{\mathrm{Im}}}
\def\Coker{\mathop{\mathrm{Coker}}}
\def\Coim{\mathop{\mathrm{Coim}}}

% Boxtimes
%
\DeclareMathSymbol{\boxtimes}{\mathbin}{AMSa}{"02}

%
% Macros for moduli stacks/spaces
%
\def\QCohstack{\mathcal{QC}\!\mathit{oh}}
\def\Cohstack{\mathcal{C}\!\mathit{oh}}
\def\Spacesstack{\mathcal{S}\!\mathit{paces}}
\def\Quotfunctor{\mathrm{Quot}}
\def\Hilbfunctor{\mathrm{Hilb}}
\def\Curvesstack{\mathcal{C}\!\mathit{urves}}
\def\Polarizedstack{\mathcal{P}\!\mathit{olarized}}
\def\Complexesstack{\mathcal{C}\!\mathit{omplexes}}
% \Pic is the operator that assigns to X its picard group, usage \Pic(X)
% \Picardstack_{X/B} denotes the Picard stack of X over B
% \Picardfunctor_{X/B} denotes the Picard functor of X over B
\def\Pic{\mathop{\mathrm{Pic}}\nolimits}
\def\Picardstack{\mathcal{P}\!\mathit{ic}}
\def\Picardfunctor{\mathrm{Pic}}
\def\Deformationcategory{\mathcal{D}\!\mathit{ef}}
\usepackage{cleveref}
\crefname{originallemma}{original lemma}{original lemmas}
\Crefname{originallemma}{Original lemma}{Original lemmas}
\makeatletter
\newlabel{lemma-descending-properties-morphisms}{{041J}{source}{}{sourcepointer-1}{}}
\newlabel{lemma-descending-properties-morphisms@cref}{{[originallemma][0][]041J}{[1][0][]source}{}{}{}}
\makeatother
\crefname{originalreference}{original reference}{original references}
\Crefname{originalreference}{Original reference}{Original references}
\makeatletter
\newlabel{spaces-morphisms-lemma-affine-local}{{03WG}{source}{}{sourcepointer-2}{}}
\newlabel{spaces-morphisms-lemma-affine-local@cref}{{[originalreference][0][]03WG}{[1][0][]source}{}{}{}}
\makeatother
\crefname{originalreference}{original reference}{original references}
\Crefname{originalreference}{Original reference}{Original references}
\makeatletter
\newlabel{descent-definition-descent-datum}{{023V}{source}{}{sourcepointer-3}{}}
\newlabel{descent-definition-descent-datum@cref}{{[originalreference][0][]023V}{[1][0][]source}{}{}{}}
\makeatother
\crefname{originalreference}{original reference}{original references}
\Crefname{originalreference}{Original reference}{Original references}
\makeatletter
\newlabel{descent-lemma-descent-data-sheaves}{{02W5}{source}{}{sourcepointer-4}{}}
\newlabel{descent-lemma-descent-data-sheaves@cref}{{[originalreference][0][]02W5}{[1][0][]source}{}{}{}}
\makeatother
\crefname{originalreference}{original reference}{original references}
\Crefname{originalreference}{Original reference}{Original references}
\makeatletter
\newlabel{descent-lemma-affine}{{0245}{source}{}{sourcepointer-5}{}}
\newlabel{descent-lemma-affine@cref}{{[originalreference][0][]0245}{[1][0][]source}{}{}{}}
\makeatother
\crefname{originalreference}{original reference}{original references}
\Crefname{originalreference}{Original reference}{Original references}
\makeatletter
\newlabel{descent-lemma-descending-property-affine}{{02L5}{source}{}{sourcepointer-6}{}}
\newlabel{descent-lemma-descending-property-affine@cref}{{[originalreference][0][]02L5}{[1][0][]source}{}{}{}}
\makeatother
\crefname{originalreference}{original reference}{original references}
\Crefname{originalreference}{Original reference}{Original references}
\makeatletter
\newlabel{spaces-properties-proposition-sheaf-fpqc}{{0APL}{source}{}{sourcepointer-7}{}}
\newlabel{spaces-properties-proposition-sheaf-fpqc@cref}{{[originalreference][0][]0APL}{[1][0][]source}{}{}{}}
\makeatother
\crefname{originallemma}{original lemma}{original lemmas}
\Crefname{originallemma}{Original lemma}{Original lemmas}
\makeatletter
\newlabel{lemma-descending-property-isomorphism}{{041Y}{source}{}{sourcepointer-8}{}}
\newlabel{lemma-descending-property-isomorphism@cref}{{[originallemma][0][]041Y}{[1][0][]source}{}{}{}}
\makeatother
\title{Being affine is fpqc local on the base}
\author{The Stacks Project Authors}
\date{stacks-project:spaces-descent.tex:local-snapshot-20261012; GNU FDL 1.2 or later}
\begin{document}
\maketitle
\noindent\textit{Editorial note.} All delivered mathematical text is the author's own, copied verbatim from the pinned Stacks Project chapter file spaces-descent.tex; no statement or proof step was rewritten, reordered or removed.\par\medskip
\noindent\textit{Editorial note.} The author's parenthetical '(verification omitted, compare with Descent, Lemma descent-lemma-descent-data-sheaves)' inside the descent-datum step is kept verbatim: it omits only the routine cocycle check of \textbackslash{}varphi, with the schemes-side comparison given, while the main line (effectivity, sheaf condition, isomorphism descent) is fully derived.\par\medskip
\noindent\textit{Editorial note.} The counted span cites lemmas in Descent, Morphisms of Spaces and Properties of Spaces plus same-chapter lemmas, recorded in assembly.external\_references with canonical tags as navigation pointers.\par\medskip
\noindent\textit{Allowed standard prerequisites (editorial).}\begin{itemize}
\item Framework lemma for descending properties of morphisms. Stacks Project, Descent on Spaces, Lemma lemma-descending-properties-morphisms. Supplies the criterion structure (assumptions (1), (2)) that the proof verifies.
\item Affineness is local in the \textbackslash{}'etale topology. Stacks Project, Morphisms of Spaces, Lemma spaces-morphisms-lemma-affine-local. Supplies assumptions (1) and (2).
\item Effectivity of affine descent data. Stacks Project, Descent, Lemma descent-lemma-affine (with Definition descent-definition-descent-datum). Supplies the scheme Y \textbackslash{}to Z affine from the descent datum.
\item Algebraic spaces satisfy the fpqc sheaf condition. Stacks Project, Properties of Spaces, Proposition spaces-properties-proposition-sheaf-fpqc. Supplies the morphism Y \textbackslash{}to X from the descent datum element.
\item fpqc descent of the isomorphism property. Stacks Project, Descent on Spaces, Lemma lemma-descending-property-isomorphism. Supplies that Y \textbackslash{}to X is an isomorphism from its base change being one.
\end{itemize}\medskip
\ifdefined\hypertarget\hypertarget{sourcepointer-1}{}\fi
\noindent\textit{Source reference.} Original lemma 041J: sources/stacks-chow/tags.txt, source lines 5294--5294. Cross-chapter standard prerequisite cited by the author inside the counted span; it is used as a pointer and its own proof is not part of this item.\par\medskip
\ifdefined\hypertarget\hypertarget{sourcepointer-2}{}\fi
\noindent\textit{Source reference.} Original reference 03WG: sources/stacks-chow/tags.txt, source lines 5111--5111. Cross-chapter standard prerequisite cited by the author inside the counted span; it is used as a pointer and its own proof is not part of this item.\par\medskip
\ifdefined\hypertarget\hypertarget{sourcepointer-3}{}\fi
\noindent\textit{Source reference.} Original reference 023V: sources/stacks-chow/tags.txt, source lines 2756--2756. Cross-chapter standard prerequisite cited by the author inside the counted span; it is used as a pointer and its own proof is not part of this item.\par\medskip
\ifdefined\hypertarget\hypertarget{sourcepointer-4}{}\fi
\noindent\textit{Source reference.} Original reference 02W5: sources/stacks-chow/tags.txt, source lines 3792--3792. Cross-chapter standard prerequisite cited by the author inside the counted span; it is used as a pointer and its own proof is not part of this item.\par\medskip
\ifdefined\hypertarget\hypertarget{sourcepointer-5}{}\fi
\noindent\textit{Source reference.} Original reference 0245: sources/stacks-chow/tags.txt, source lines 2766--2766. Cross-chapter standard prerequisite cited by the author inside the counted span; it is used as a pointer and its own proof is not part of this item.\par\medskip
\ifdefined\hypertarget\hypertarget{sourcepointer-6}{}\fi
\noindent\textit{Source reference.} Original reference 02L5: sources/stacks-chow/tags.txt, source lines 3396--3396. Cross-chapter standard prerequisite cited by the author inside the counted span; it is used as a pointer and its own proof is not part of this item.\par\medskip
\ifdefined\hypertarget\hypertarget{sourcepointer-7}{}\fi
\noindent\textit{Source reference.} Original reference 0APL: sources/stacks-chow/tags.txt, source lines 13479--13479. Cross-chapter standard prerequisite cited by the author inside the counted span; it is used as a pointer and its own proof is not part of this item.\par\medskip
\ifdefined\hypertarget\hypertarget{sourcepointer-8}{}\fi
\noindent\textit{Source reference.} Original lemma 041Y: sources/stacks-chow/tags.txt, source lines 5309--5309. Cross-chapter standard prerequisite cited by the author inside the counted span; it is used as a pointer and its own proof is not part of this item.\par\medskip
\begin{lemma}
\label{lemma-descending-property-affine}
The property $\mathcal{P}(f) =$``$f$ is affine''
is fpqc local on the base.
\end{lemma}

\begin{proof}
We will use
Lemma \ref{lemma-descending-properties-morphisms}
to prove this. Assumptions (1) and (2) of that lemma follow from
Morphisms of Spaces,
Lemma \ref{spaces-morphisms-lemma-affine-local}.
Let $Z' \to Z$ be a surjective flat morphism of affine schemes over $S$.
Let $f : X \to Z$ be a morphism of algebraic spaces, and assume
that the base change $f' : Z' \times_Z X \to Z'$ is affine.
Let $X'$ be a scheme representing $Z' \times_Z X$.
We obtain a canonical isomorphism
$$
\varphi : X' \times_Z Z' \longrightarrow Z' \times_Z X'
$$
since both schemes represent the algebraic space $Z' \times_Z Z' \times_Z X$.
This is a descent datum for $X'/Z'/Z$, see
Descent, Definition \ref{descent-definition-descent-datum}
(verification omitted, compare with
Descent, Lemma \ref{descent-lemma-descent-data-sheaves}).
Since $X' \to Z'$ is affine this descent datum is effective, see
Descent, Lemma \ref{descent-lemma-affine}.
Thus there exists a scheme $Y \to Z$ over $Z$ and an
isomorphism $\psi : Z' \times_Z Y \to X'$ compatible with descent data.
Of course $Y \to Z$ is affine (by construction or by
Descent, Lemma \ref{descent-lemma-descending-property-affine}).
Note that $\mathcal{Y} = \{Z' \times_Z Y \to Y\}$ is a
fpqc covering, and interpreting $\psi$ as an element of
$X(Z' \times_Z Y)$ we see that $\psi \in \check{H}^0(\mathcal{Y}, X)$.
By the sheaf condition for $X$ with respect to this covering (see
Properties of Spaces, Proposition
\ref{spaces-properties-proposition-sheaf-fpqc})
we obtain a morphism $Y \to X$.
By construction the base change of this to $Z'$ is an isomorphism, hence
an isomorphism by
Lemma \ref{lemma-descending-property-isomorphism}.
This proves that $X$ is representable by an affine scheme and we win.
\end{proof}

\end{document}
